\documentclass[11pt]{article}

% ============================================================
%  The Delta F = 2 framework and the RS flavour/CP problem
%  A slow, pedagogical internal review.
%  Written to track (a) the physics, (b) what is in OUR code,
%  (c) the subtleties and gaps.
%  This is the shared deep-dive that the four meson-mixing
%  notes (epsilon_K, D0, Delta m_d, Delta m_s) build on, and
%  the authoritative normalization reference for them.
%  Internal working note, not for distribution.
% ============================================================

\usepackage[margin=1in]{geometry}
\usepackage{amsmath,amssymb,amsthm}
\usepackage{mathtools}
\usepackage{graphicx}
\usepackage{booktabs}
\usepackage{enumitem}
\usepackage{xcolor}
\usepackage{framed}
\usepackage[colorlinks=true,linkcolor=blue,citecolor=blue,urlcolor=blue]{hyperref}

% lightweight "callout" box
\colorlet{shadecolor}{gray!12}
\newenvironment{callout}[1]{%
  \par\medskip\begin{shaded}\noindent\textbf{#1}\par\smallskip}{\end{shaded}\medskip}

\newcommand{\MKK}{M_{\mathrm{KK}}}
\newcommand{\LIR}{\Lambda_{\mathrm{IR}}}
\newcommand{\half}{\tfrac{1}{2}}
\newcommand{\vev}{v}
\newcommand{\eps}{\varepsilon}
\newcommand{\Mtwelve}{M_{12}}
\newcommand{\ImMtwelve}{\mathrm{Im}\,M_{12}}
\newcommand{\gss}{g_{s*}}
\newcommand{\Nc}{N_c}
\newcommand{\QVLL}{Q_1^{\mathrm{VLL}}}
\newcommand{\QVRR}{Q_1^{\mathrm{VRR}}}
\newcommand{\QLRfour}{Q_4^{\mathrm{LR}}}
\newcommand{\QLRfive}{Q_5^{\mathrm{LR}}}
\newcommand{\CVLL}{C_1^{\mathrm{VLL}}}
\newcommand{\CVRR}{C_1^{\mathrm{VRR}}}
\newcommand{\CLRfour}{C_4^{\mathrm{LR}}}
\newcommand{\CLRfive}{C_5^{\mathrm{LR}}}
\newcommand{\Heff}{\mathcal{H}_{\mathrm{eff}}}
\newcommand{\TA}{T^A}
\newcommand{\code}[1]{\texttt{#1}}
\newcommand{\todo}[1]{{\color{red}[#1]}}

\theoremstyle{definition}
\newtheorem{definition}{Definition}

\title{\bf The $\Delta F=2$ Framework and the RS Flavour/CP Problem\\[4pt]
\large A slow, pedagogical internal review --- physics, code, and the
authoritative normalization reference}
\author{Internal working note (5D--Neutrino--Mixing repo)}
\date{\today}

\begin{document}
\maketitle

\begin{abstract}
\noindent
This note is the shared, from-scratch deep-dive that the four meson-mixing
constraints in our catalogue --- $\eps_K$, $D^0$--$\bar D^0$, $\Delta m_d$,
$\Delta m_s$ --- all stand on. It assumes no prior knowledge of effective
Hamiltonians, four-quark operators, or the ``RS flavour problem,'' and builds
each from the ground up. Part~\ref{part:framework} constructs the $\Delta F=2$
effective Hamiltonian and the operator basis that any neutral-meson mixing
analysis uses --- the vector left-left operator $\QVLL$ and the
chirality-flipping left--right operators $\QLRfour,\QLRfive$ --- and explains
why the LR operators carry a large chiral and QCD-running enhancement, making
them the dominant and dangerous contribution. Part~\ref{part:rs} shows how a
warped extra dimension turns these operators on at tree level through
Kaluza--Klein (KK) gluon exchange, how the RS-GIM/anarchy structure feeds in,
and states the unified RS flavour/CP problem ($\eps_K$ worst at $\sim20$~TeV
anarchic, the heavy-meson and charm bounds milder, custodial symmetry
\emph{useless} here). Part~\ref{part:code} documents our implementation
(\code{quarkConstraints/deltaf2.py}) line by line, including the load-bearing
post-B3 M12-ready matrix elements and the unchanged $(1/6)$ Wilson colour factor.
Part~\ref{part:subtle}
records the subtleties, the deferred pieces, and the documented hazards. An
annotated reading guide closes the note. This note is the \emph{authoritative
normalization reference} for the four meson notes: when in doubt about a factor,
read \S\ref{sec:r5} here.
\end{abstract}

\tableofcontents

% ====================================================================
\part*{Orientation: why this note exists}
\addcontentsline{toc}{section}{Orientation: why this note exists}
% ====================================================================

Our constraint catalogue contains four neutral-meson mixing observables. They
look superficially independent --- a kaon CP parameter, a charm mass splitting,
two beauty mass splittings --- but they are not. All four are the \emph{same}
short-distance physics: a $\Delta F=2$ (two units of flavour change)
four-quark amplitude generated, in a warped model, by tree-level KK-gluon
exchange. They share one effective Hamiltonian, one operator basis, one matching
calculation, and --- in our code --- one core module,
\code{quarkConstraints/deltaf2.py}. Writing the same derivation four times would
be wasteful and, worse, error-prone: a normalization choice made differently in
two of the notes would be a silent bug. So we factor the common physics out into
this single note, derive everything once, pin the normalization convention once,
and let the four observable-specific notes specialise.

There is a second, sharper reason this note exists. The single most consequential
subtlety in the whole quark-flavour pipeline is a normalization convention: the
Wilson coefficient carries a $1/6$ colour factor, while the M12-ready matrix
element carries the post-B3 $\tfrac13$ prefactor. The older pre-B3 review text
mixed this with a stale $\tfrac23$ M12-ready value and a swapped, doubled LR pair.
That is why this note is tagged the authoritative reference: any future edit to a
meson note that touches a factor must be checked against \S\ref{sec:r5}.

A one-paragraph executive summary, to be unpacked over the following pages:

\begin{quote}
\small
Below the KK scale, integrating out a heavy KK gluon leaves a set of local
four-quark operators times Wilson coefficients --- the $\Delta F=2$ effective
Hamiltonian. A KK gluon couples to \emph{both} chiralities, so it generates not
only the vector left-left operator that the Standard Model (SM) box already makes,
but also chirality-flipping left--right operators. The LR operators are
\emph{doubly enhanced}: by a chiral factor at the hadronic level (the
pseudoscalar density is large for a light meson) and by QCD running between the
multi-TeV matching scale and the few-GeV hadronic scale. In an \emph{anarchic}
warped model the off-diagonal KK-gluon couplings are $\mathcal{O}(1)$ in phase
and suppressed only by the same small wavefunction overlaps that set the quark
masses (the \emph{RS-GIM} mechanism). That partial suppression is not enough for
the chirally- and RG-enhanced LR operator carrying a generic CP phase: the
CP-violating part of kaon mixing, $\eps_K$, then forces the lightest KK gluon
above $\sim20$~TeV. The beauty and charm bounds are milder. Custodial symmetry,
which rescues the electroweak sector, does \emph{nothing} for these
KK-\emph{gluon} amplitudes; what helps is flavour alignment, larger bulk masses,
or a bulk flavour symmetry.
\end{quote}

\noindent The logical spine of the whole note, in one picture:
\begin{footnotesize}
\begin{verbatim}
  heavy KK gluon (mass M_KK)  --integrate out-->  H_eff^{dF=2} = sum_i C_i(mu) Q_i(mu)
                           |
  KK gluon couples to BOTH chiralities (g_L and g_R)
                           |
       +-- g_L x g_L  --> Q1_VLL  (vector, SM-like)          [mild]
       +-- g_R x g_R  --> Q1_VRR  (parity image of VLL)      [mild]
       +-- g_L x g_R  --> Q4_LR, Q5_LR  (chirality-FLIP)     [DANGEROUS]
                           |
  why dangerous:  <Q4_LR> / <Q1_VLL>  ~  chiral factor r_chi  x  QCD running
                                       ~     26 (kaon)        x     ~8        ~ 100x
                           |
  RS-GIM:  off-diag coupling ~ g_s* f_i f_j   (same f's as quark masses)  --PARTIAL
           but surviving residue is O(1) COMPLEX --> unprotected CP phase
                           |
  ==> eps_K (Im part) ~ M_KK floor ~ 20 TeV (anarchic)  --  THE RS flavour/CP problem
      Delta m_d / Delta m_s / D  (real part, heavy mesons)  --  MILDER (no big chiral factor)
                           |
  custodial SU(2)_R x P_LR  fixes T and Z b_L b_L  ==  does NOTHING for KK-gluon dF=2
  what helps: ALIGNMENT / larger bulk masses / bulk flavour symmetry
\end{verbatim}
\end{footnotesize}
\noindent Keep this map in mind; every box is unpacked below. The most common
confusion --- that ``custodial symmetry is the cure for everything in RS'' --- is
already defused here: the custodial arrow lands on the electroweak boxes, not on
the $\Delta F=2$ gluon box.

% ====================================================================
\part{The $\Delta F=2$ framework, from scratch}
\label{part:framework}
% ====================================================================

\section{Mixing, $M_{12}$, and what ``$\Delta F=2$'' means}
\label{sec:mixing}

A neutral meson $P^0$ (here $K^0=\bar s d$, $D^0=c\bar u$, $B_d^0=\bar b d$, or
$B_s^0=\bar b s$) and its antiparticle $\bar P^0$ carry opposite values of a
flavour quantum number $F$ (strangeness, charm, or beauty). Any interaction that
turns $P^0$ into $\bar P^0$ changes $F$ by two units --- this is what
``\emph{$\Delta F=2$}'' names. Such an interaction makes the flavour eigenstates
$P^0,\bar P^0$ \emph{not} the mass eigenstates: the physical states are
superpositions, and as a beam of $P^0$ propagates it oscillates into $\bar P^0$
and back. That oscillation is \emph{meson mixing}.

The $2\times2$ effective Hamiltonian governing the $P^0$--$\bar P^0$ system splits
into a dispersive (off-shell) part $\Mtwelve$ and an absorptive (on-shell) part
$\Gamma_{12}$. The single quantity that all four of our observables care about is
the off-diagonal dispersive element $\Mtwelve$, because the mass splitting and the
CP-violating phase of the mixing are
\begin{equation}
\Delta m_P \;\simeq\; 2\,|\Mtwelve|, \qquad
\phi_P \;=\; \arg\Mtwelve ,
\label{eq:dmphi}
\end{equation}
up to small absorptive corrections. New physics enters as an additive shift
$\Mtwelve = \Mtwelve^{\rm SM} + \Mtwelve^{\rm NP}$, and our entire job is to
compute the warped-model $\Mtwelve^{\rm NP}$ and ask whether it overshoots the
room left by data. For the kaon there is a special twist: the binding quantity is
not the magnitude but the \emph{imaginary} part, the indirect-CP parameter
\begin{equation}
\eps_K \;\simeq\; \frac{\kappa_\eps}{\sqrt2\,\Delta m_K}\,\ImMtwelve ,
\label{eq:epsk}
\end{equation}
which is why $\eps_K$ is, as we will see, the most constraining of the four.

\section{The effective Hamiltonian: integrating out the heavy state}
\label{sec:effham}

The new physics that mixes $P^0$ and $\bar P^0$ lives at a scale $\MKK$ far above
the meson mass. We never need its full dynamics; we need only its low-energy
\emph{shadow}. The standard tool is the operator product expansion: replace the
exchange of the heavy state by a tower of local (contact) operators built from
the light quark fields, each multiplied by a coefficient that encodes the
high-scale physics. The result is the $\Delta F=2$ effective Hamiltonian
\begin{equation}
\boxed{\;\Heff^{\Delta F=2} \;=\; \sum_i C_i(\mu)\, Q_i(\mu).\;}
\label{eq:heff}
\end{equation}
The split is the whole point. The \textbf{Wilson coefficients} $C_i(\mu)$ are
numbers: they carry the short-distance physics (the heavy propagator, the
couplings, the matching), and depend on the renormalization scale $\mu$. The
\textbf{operators} $Q_i(\mu)$ are local products of four quark fields: their
hadronic matrix elements $\langle\bar P^0|Q_i|P^0\rangle$ carry the long-distance,
non-perturbative physics (how the quarks bind into the meson), and also depend on
$\mu$. The $\mu$-dependence cancels between the two in any physical amplitude.
Schematically,
\begin{equation}
\Mtwelve^{\rm NP} \;=\; \sum_i C_i(\mu)\,\langle\bar P^0|Q_i|P^0\rangle(\mu).
\label{eq:m12sum}
\end{equation}
The art of a $\Delta F=2$ computation is (i) \emph{matching}: finding the $C_i$ at
the high scale $\mu=\MKK$ by comparing the full heavy-state amplitude to the
effective one; (ii) \emph{running}: evolving the $C_i$ down to a low scale
$\mu\approx2$~GeV with the QCD renormalization group; and (iii) \emph{hadronic
input}: supplying the matrix elements $\langle Q_i\rangle$ at that scale from
lattice QCD. Our code does all three; the rest of this part sets up the
operators those three steps act on.

\section{The operator basis and its chiral structure}
\label{sec:basis}

The full $\Delta F=2$ basis has several Lorentz and colour structures, but for a
KK-gluon source only three matter. Write $P_{L,R}=\half(1\mp\gamma_5)$ for the
chirality projectors and $\alpha,\beta$ for colour indices. The three operators
are
\begin{align}
\QVLL  &= (\bar q_\alpha \gamma_\mu P_L q'_\alpha)(\bar q_\beta \gamma^\mu P_L q'_\beta),
\label{eq:QVLL}\\
\QLRfour &= (\bar q_\alpha P_L q'_\alpha)(\bar q_\beta P_R q'_\beta),
\label{eq:QLR4}\\
\QLRfive &= (\bar q_\alpha P_L q'_\beta)(\bar q_\beta P_R q'_\alpha),
\label{eq:QLR5}
\end{align}
where $q,q'$ are the two quark flavours of the meson (e.g.\ $q=s,\ q'=d$ for the
kaon). There is also a parity image $\QVRR$, obtained from $\QVLL$ by $L\to R$;
its matrix element equals that of $\QVLL$ by parity, and our code uses exactly
that fact. The names are standard:
\begin{itemize}[leftmargin=1.4em]
\item \textbf{VLL} (``vector, left, left''): a vector \emph{current} of
left-handed fields, squared. Both bilinears are $\bar q\gamma_\mu P_L q'$. This is
the only structure the SM box diagram makes, because the $W$ boson couples only to
left-handed quarks.
\item \textbf{LR} (``left--right''): a product of a left-handed and a
right-handed \emph{scalar} (density) bilinear, $\bar q P_L q'$ times
$\bar q P_R q'$. These are \emph{chirality-flipping}: they require the new physics
to couple to both chiralities. A KK gluon, like the ordinary gluon, is
vector-like and couples to left- and right-handed quarks alike --- so it makes
LR operators.
\end{itemize}
$\QLRfour$ and $\QLRfive$ differ only in their colour contraction ($\QLRfour$
colour-singlet $\times$ colour-singlet, $\QLRfive$ colour-crossed). The decisive
physics is that the LR operators are \emph{hadronically enhanced} relative to
VLL, and we turn to that now.

\section{Why the LR operators dominate: the chiral enhancement}
\label{sec:chiral}

The matrix element of a left--right \emph{scalar} operator is larger than that of
a left-left \emph{vector} operator by the ``chiral factor''
\begin{equation}
r_\chi \;\equiv\; \left(\frac{m_P}{m_q(\mu)+m_{q'}(\mu)}\right)^2 .
\label{eq:rchi}
\end{equation}
The origin is the partial-conservation argument: the divergence of the
axial current relates the pseudoscalar density $\bar q\gamma_5 q'$ to
$m_P^2/(m_q+m_{q'})$ via the equations of motion. For a \emph{light} meson built
of \emph{light} quarks this denominator is tiny and $r_\chi$ is large. The numbers
are the entire story of why $\eps_K$ is the worst case --- here computed with the
exact light-quark masses our code feeds the matrix elements:
\begin{center}
\begin{tabular}{lcccc}
\toprule
System & $m_P$ (GeV) & $m_q+m_{q'}$ (GeV) & $r_\chi$ & comment \\
\midrule
$K^0$ ($\bar s d$)  & $0.498$ & $0.0934+0.00467$ & $\approx 26$ & light meson, huge $r_\chi$ \\
$D^0$ ($c\bar u$)   & $1.865$ & $1.27+0.00216$   & $\approx 2.2$ & charm quark in the denominator \\
$B_d^0$ ($\bar b d$)& $5.280$ & $4.18+0.00467$   & $\approx 1.6$ & $b$ quark dominates \\
$B_s^0$ ($\bar b s$)& $5.367$ & $4.18+0.0934$    & $\approx 1.6$ & $b$ quark dominates \\
\bottomrule
\end{tabular}
\end{center}
The lesson is stark. For the kaon the LR operator is chirally enhanced by a factor
$\sim26$; for the heavy mesons by less than $\sim2$, because the heavy quark in
the denominator of \eqref{eq:rchi} kills the enhancement. \emph{This single table
is the deepest reason $\eps_K$ pushes $\MKK$ an order of magnitude higher than
$\Delta m_{d,s}$ or $D$ mixing do.} The chiral factor is the first of two LR
enhancements; the second is QCD running.

\section{The second enhancement: QCD running}
\label{sec:running}

The Wilson coefficients are matched at $\MKK$ (multi-TeV) but the matrix elements
are evaluated at $\mu_{\rm had}\approx2$~GeV. Between those scales QCD evolves the
coefficients, and the effect is not small:
\begin{itemize}[leftmargin=1.4em]
\item the vector operators $\QVLL,\QVRR$ run multiplicatively with a single
anomalous dimension and change only modestly;
\item the scalar LR pair $\QLRfour,\QLRfive$ \emph{mix} under running through a
$2\times2$ anomalous-dimension matrix, and the running \emph{amplifies} the LR
contribution by a factor often quoted as $\eta_1^{-5}\sim8$ in the literature.
\end{itemize}
Combined with the chiral factor of \S\ref{sec:chiral}, the running brings the
total kaon LR enhancement to roughly \emph{two orders of magnitude} over the naive
VLL estimate: the chiral factor ($r_\chi\approx26$ with the repo's $2$~GeV
light-quark masses) times the running ($\eta_1^{-5}\sim8$) lands near $\sim200$
before bag/coefficient details are folded in. The split is worth stating cleanly
because the two enhancements are independent and have different origins --- the
chiral factor is hadronic (it lives in the matrix element, \S\ref{sec:chiral}),
the $\eta_1^{-5}$ is renormalization-group (it lives in the Wilson coefficient).
Note also that a \emph{lighter} light-quark input would push $r_\chi$
\emph{higher}, not lower (the mass sits in the denominator of \eqref{eq:rchi}), so
the repo's relatively heavy $2$~GeV masses are the conservative choice here. We
document the exact anomalous-dimension constants our code uses
(\code{\_GAMMA\_VLL = 4.0} for the vector sector and a $2\times2$
\code{\_GAMMA\_LR} matrix for the LR sector) in \S\ref{sec:codeRunning}.

\begin{callout}{Common confusions, defused.}
\begin{itemize}[leftmargin=1.2em,nosep]
\item[\textbf{A.}] \emph{``The KK gluon makes only the SM operator.''} No --- a
KK gluon is vector-like and couples to both chiralities, so beyond $\QVLL$ it
generates the chirality-flipping $\QLRfour,\QLRfive$, which are absent in the SM
box and are the \emph{dominant} contribution. (\S\ref{sec:basis})
\item[\textbf{B.}] \emph{``LR dominates because its coefficient is large.''} No
--- its coefficient is comparable to VLL's. LR dominates because its \emph{matrix
element} is enhanced by the chiral factor $r_\chi$ and by QCD running ---
hadronic and RG physics, not the short-distance coefficient.
(\S\ref{sec:chiral}--\ref{sec:running})
\item[\textbf{C.}] \emph{``All four mesons are equally constrained.''} No ---
$r_\chi\approx26$ for the kaon but $\lesssim2$ for the heavy mesons, so the LR
enhancement, and hence the $\MKK$ floor, is far larger for $\eps_K$.
(\S\ref{sec:chiral})
\item[\textbf{D.}] \emph{``$\Delta m_K$ and $\eps_K$ are the same constraint.''}
No --- $\Delta m_K\simeq2\,\mathrm{Re}\,\Mtwelve$ is the magnitude;
$\eps_K\propto\ImMtwelve$ is the CP-violating phase, and is roughly an order of
magnitude tighter. (\S\ref{sec:mixing})
\end{itemize}
\end{callout}

% ====================================================================
\part{How warped models generate $\Delta F=2$}
\label{part:rs}
% ====================================================================

\section{RS recap and the seed of flavour violation}
\label{sec:rsrecap}

Spacetime is a slice of AdS$_5$ between a UV brane (Planck scale) and an IR brane
(TeV scale); the warp factor is $\eps = e^{-k\pi r_c}\sim10^{-15}$. A bulk fermion
carries a dimensionless mass parameter $c$ that fixes where its massless zero mode
lives: $c>\half$ peaks toward the UV brane (small IR overlap, a light SM fermion),
$c<\half$ peaks toward the IR brane (large overlap, a heavy fermion). The
normalised value of the zero-mode wavefunction at the IR brane is the overlap
factor $f(c)$, and an $\mathcal{O}(1)$ spread of $c$ parameters reproduces the SM
mass hierarchy with no hierarchical 5D Yukawas --- the central virtue of the
construction.

The price is flavour. A bulk gauge field (gluon, $W$, $Z$, photon) has a tower of
KK modes localised toward the IR brane. The first KK gluon couples to a quark with
strength set by the \emph{overlap} of the quark's profile with the IR-peaked KK
profile: light (UV-localised) quarks couple weakly, heavy (IR-localised) quarks
strongly. In the quark mass basis the KK-gluon coupling matrix is therefore
\emph{non-universal} and \emph{not} flavour-diagonal: rotating to the mass basis
leaves off-diagonal entries. Schematically the off-diagonal KK-gluon coupling
between the $i,j$ generations is
\begin{equation}
(g_{L})_{ij} \;\sim\; \gss\, f_{q_i} f_{q_j}, \qquad
(g_{R})_{ij} \;\sim\; \gss\, f_{q'_i} f_{q'_j},
\label{eq:gij}
\end{equation}
with $\gss$ the (volume-enhanced) 5D strong coupling and $f$ the same overlap
factors that set the masses. This is the seed of every $\Delta F=2$ problem in RS.

\section{Tree-level KK-gluon exchange and the matching}
\label{sec:treelevel}

Integrate out the first KK gluon. Two quark currents exchange it in the
$t$-channel; the heavy propagator gives a $1/\MKK^2$, and the two vertices give
two powers of the coupling \eqref{eq:gij}. A KK gluon, being a colour octet,
carries an $SU(3)$ generator $\TA$ at each vertex, so the colour structure of the
amplitude is $\TA\otimes\TA$. The exact $SU(3)$ Fierz identity rearranges this
into colour singlets,
\begin{equation}
\sum_A (\TA)_{\alpha\beta}(\TA)_{\gamma\delta}
= \half\left(\delta_{\alpha\delta}\delta_{\gamma\beta}
- \frac{1}{\Nc}\,\delta_{\alpha\beta}\delta_{\gamma\delta}\right),
\qquad \Nc=3,
\label{eq:fierz}
\end{equation}
which is what brings the result into the operator basis of \S\ref{sec:basis}:
\begin{itemize}[leftmargin=1.4em]
\item a left current $\times$ a left current gives $\QVLL$, with coefficient
$\propto (g_L)_{ij}^2/\MKK^2$;
\item a right $\times$ right gives $\QVRR$, $\propto (g_R)_{ij}^2/\MKK^2$;
\item a left $\times$ right gives the chirality-flipping $\QLRfour,\QLRfive$,
$\propto (g_L)_{ij}(g_R)_{ij}/\MKK^2$.
\end{itemize}
This structure is exactly that of Csaki, Falkowski and Weiler (0804.1954, their
Eqs.~3.3--3.7), and our repo reproduces it with definite numerical colour/Fierz
factors documented line-by-line in \S\ref{sec:codeWilsons}. The essential point is
parametric: \emph{every} coefficient scales as $1/\MKK^2$ (tree-level decoupling),
and the LR coefficient is the product of a left and a right coupling, so it is
fed by both chiralities of the misalignment.

\section{The RS-GIM mechanism: partial, not complete, protection}
\label{sec:rsgim}

Equation~\eqref{eq:gij} contains a built-in suppression, and it pays to be precise
about its reach. The off-diagonal coupling is proportional to $f_{q_i}f_{q_j}$ ---
the \emph{same} small overlap factors that make the light-quark masses small. A
$\Delta F=2$ amplitude between the first two generations is therefore
automatically suppressed by the light-quark overlaps. This is the \textbf{RS-GIM
mechanism}: flavour violation involving light quarks is suppressed by the same
geometry that makes them light, mimicking the SM's own GIM suppression. It is
genuine and automatic; without it RS would be hopeless.

But it is only \emph{partial}. Three reasons it fails to rescue the dangerous
cases:
\begin{itemize}[leftmargin=1.4em]
\item \textbf{The surviving coefficients are $\mathcal{O}(1)$ complex.} What
escapes the overlap suppression carries generic, unsuppressed CP phases. Nothing
forces the surviving LR coefficient to be real or aligned with the SM, so its
imaginary part is unprotected --- fatal for $\eps_K$.
\item \textbf{The LR operator is doubly enhanced.} Even a modest
$\CLRfour$ is magnified by the chiral factor (\S\ref{sec:chiral}) and by QCD
running (\S\ref{sec:running}), giving back much of what RS-GIM suppressed ---
again worst for the kaon.
\item \textbf{The strong coupling is volume-enhanced.} The 5D gluon coupling
$\gss$ in \eqref{eq:gij} is larger than the perturbative $g_s$, partially undoing
the overlap suppression.
\end{itemize}
The net result is a residual, complex LR coefficient that, once dressed by the
enhancements, saturates the experimental constraint for KK scales an order of
magnitude above the TeV we would like.

\section{The unified RS flavour/CP problem}
\label{sec:unified}

We can now state the whole problem in one place, ranked. The off-diagonal
KK-gluon couplings are anarchic and complex; the LR operator they feed is
chirally and RG enhanced; the amplitude scales as $1/\MKK^2$; and a constraint on
the amplitude becomes a floor on $\MKK$. The floor depends on \emph{which}
observable and \emph{which} part (real vs.\ imaginary):
\begin{center}
\begin{tabular}{llll}
\toprule
Observable & Quantity & Why this severity & $\MKK$ floor (anarchic) \\
\midrule
$\eps_K$ & $\ImMtwelve$ (CP) & big $r_\chi\approx26$ + no SM Im to hide behind & $\sim 20$~TeV \emph{(worst)} \\
$\Delta m_s$ & $|\Mtwelve|$ & small $r_\chi\approx1.6$; clean SM & $\sim 2$--$3$~TeV \\
$\Delta m_d$ & $|\Mtwelve|$ & small $r_\chi\approx1.6$; clean SM & $\sim 2$--$3$~TeV \\
$D^0$ & $|\Mtwelve|$ & up-sector; SM long-distance dirty $\Rightarrow$ loose budget & $\sim$few~TeV \\
\bottomrule
\end{tabular}
\end{center}
The hierarchy of floors is entirely the hierarchy of chiral factors in
\S\ref{sec:chiral} (modulated by how clean each SM prediction is). The kaon's
CP-violating channel is the binding edge --- this is the celebrated
\textbf{RS flavour/CP problem}: with anarchic $\mathcal{O}(1)$ complex 5D Yukawas,
$\eps_K$ forces the lightest KK gluon to $\sim20$~TeV
(Csaki--Falkowski--Weiler quote $M_G\gtrsim22\pm6$~TeV from $\eps_K$, rising to
$\sim33$~TeV in pseudo-Goldstone composite-Higgs variants). The asymmetry between
the CP and CP-conserving parts is large: in their analysis the bound on the
\emph{imaginary} part of the LR coefficient is $\sim1.6\times10^5$~TeV against
$\sim1.2\times10^4$~TeV from the \emph{real} part --- a factor $\sim13$ tighter on
the coefficient, which is exactly why $\eps_K$ dominates $\Delta m_K$.

\begin{callout}{Worked example --- the LR enhancement chain for a $3$~TeV kaon point (schematic).}
We trace the chain of magnitudes with round numbers; the code does each step
exactly (\S\ref{part:code}). Take $\MKK=3$~TeV and a typical anarchic, complex
$s$--$d$ coupling. First, the \emph{hadronic} enhancement of the LR operator over
the vector operator, using the repo matrix elements
(\S\ref{sec:codeME}) with $f_K=0.1557$~GeV, $m_K=0.498$~GeV, $r_\chi\approx26$:
\[
\frac{\langle\QLRfour\rangle}{\langle\QVLL\rangle}
= \frac{\big(\tfrac14 r_\chi+\tfrac{1}{24}\big)B_4^K}{\tfrac13 B_1^K}
\approx 32 ,
\qquad\Big(\text{and }\frac{\langle\QLRfive\rangle}{\langle\QVLL\rangle}\approx8.6\Big),
\]
then QCD running multiplies the LR piece by a factor of a few (\S\ref{sec:running}):
\[
32\times 3.5 \;\approx\; \sim 10^2 .
\]
So the LR operator, although born with a Wilson coefficient comparable to VLL,
dominates $\Mtwelve^{\rm NP}$ by roughly two orders of magnitude. Feeding the
\emph{imaginary} part into the master formula \eqref{eq:epsk} with
$\Delta m_K=3.48\times10^{-15}$~GeV and $\kappa_\eps=0.94$ produces an
$|\eps_K^{\rm NP}|$ that, for an $\mathcal{O}(1)$ phase at $3$~TeV, overshoots the
new-physics budget (central residual $\approx6.7\times10^{-5}$) by a large factor.
Since the amplitude is $\propto1/\MKK^2$, suppressing it below the budget requires
raising $\MKK$ by $\sim\sqrt{\text{overshoot}}$ --- about an order of magnitude, to
$\sim20$~TeV. \emph{This single chain --- chiral $\times$ running $\times$
$1/\Delta m_K$ amplifier on the imaginary part --- is the RS flavour/CP problem.}
For the heavy mesons the very same chain runs, but with $r_\chi\approx1.6$ the LR
enhancement is $\sim10\times$ smaller and the floor drops to a few TeV.
\end{callout}

\section{Model dependence: what custodial symmetry does \emph{not} fix}
\label{sec:custodial}

It is tempting --- and wrong --- to think that ``the custodial RS model'' cures the
problems of minimal RS across the board. Custodial symmetry
($SU(2)_L\times SU(2)_R\times U(1)_X$ in the bulk, with a discrete left--right
parity $P_{LR}$) is a genuine and important cure, but for a \emph{different} set of
problems:
\begin{itemize}[leftmargin=1.4em]
\item the continuous custodial $SU(2)_R$ removes the volume-enhanced
\emph{electroweak $T$ parameter};
\item the discrete $P_{LR}$ protects the \emph{$Z\bar b_L b_L$ vertex}.
\end{itemize}
Both are statements about \emph{electroweak} ($\Delta F=0$ or $\Delta F=1$)
physics --- gauge-boson self-energies and a flavour-diagonal $Z$ coupling. The
$\Delta F=2$ constraints here are four-quark amplitudes generated by
KK-\emph{gluon} exchange. Custodial $SU(2)_R$ and $P_{LR}$ have nothing to say
about the colour sector or about the misalignment between the gluon-coupling basis
and the quark mass basis. Both Csaki--Falkowski--Weiler (0804.1954) and
Blanke--Buras--Duling--Gori--Weiler (0809.1073) state this explicitly: the
KK-gluon $\Delta F=2$ contribution --- and the $\sim20$~TeV floor it imposes ---
survives in the custodial model essentially unchanged. The custodial construction
lowers the \emph{electroweak} floor from $\sim10$~TeV to $\sim3$~TeV, but leaves
the \emph{flavour} floor where it was.

So what \emph{does} help? Three routes, all acting on the \emph{flavour} structure
rather than the gauge structure:
\begin{enumerate}[leftmargin=1.6em]
\item \textbf{Alignment.} If the down-quark Yukawa and bulk-mass matrices are
(partially) aligned --- e.g.\ by a bulk flavour symmetry making them
simultaneously nearly diagonal --- the rotation to the mass basis leaves little
off-diagonal KK-gluon coupling, shrinking the dangerous $(g_L)_{ij}$ in
\eqref{eq:gij}. This is the warped analogue of minimal flavour violation, and it
attacks the source directly. (Compare 0907.0474, which aligns the bulk masses to
relax exactly this bound.)
\item \textbf{Larger bulk masses.} Pushing the light-generation $c$ parameters
deeper into the UV (larger $c$) suppresses the overlap factors $f_{q_i}$ and hence
the coupling \eqref{eq:gij}. This costs some naturalness of the fermion-mass
explanation but is a legitimate dial.
\item \textbf{Flavour symmetries.} Imposing a bulk $U(3)$ or $U(2)$ family
symmetry (broken in a controlled way) enforces degeneracies that suppress the
misalignment, again at the level of the couplings, not the gauge bosons.
\end{enumerate}
The common feature: all three reduce the actual quantity these constraints bound
--- $\mathrm{Im}\,(g_L)_{ij}(g_R)_{ij}$ for $\eps_K$, $|(g_L)_{ij}(g_R)_{ij}|$ for
the heavy mesons. None of them is custodial symmetry.

% ====================================================================
\part{What is in OUR code}
\label{part:code}
% ====================================================================

\emph{Reader here for the physics may skim Part~\ref{part:code} on a first pass
and jump to the subtleties in Part~\ref{part:subtle}; this part documents the
exact code so the note doubles as an implementation reference --- and the
authoritative normalization reference for the four meson notes.}

\section{Pipeline overview}
\label{sec:pipeline}

The $\Delta F=2$ physics is implemented once, in
\code{quarkConstraints/deltaf2.py}, and reached by all four meson constraints
through a thin adapter:
\begin{center}
\begin{tabular}{ll}
\toprule
Layer & File / object \\
\midrule
catalogue constraints & \code{K001} ($\eps_K$), \code{B001/B003} ($\Delta m_{d}/\Delta m_s$), \code{C001} ($D^0$) \\
import boundary (adapter) & \code{flavor\_catalog\_constraints/physics\_adapters/deltaf2.py} \\
$\Delta F=2$ core & \code{quarkConstraints/deltaf2.py} \\
QCD running & \code{quarkConstraints/qcd\_running.py} \\
\bottomrule
\end{tabular}
\end{center}
The adapter re-exports the core dataclasses unchanged and adds append-only
wrappers (e.g.\ \code{epsilon\_k\_from\_wilsons\_with\_running},
\code{bd\_mixing\_from\_wilsons\_with\_running},
\code{bs\_mixing\_from\_wilsons\_with\_running},
\code{d0\_mixing\_from\_wilsons\_with\_running}); each constraint touches the core
only through these. The model and input labels are fixed strings:
\code{DELTA\_F2\_MODEL\_V1 = "kk\_gluon\_tree\_v1"} and
\code{DELTA\_F2\_INPUT\_BUNDLE\_V1 = "deltaf2\_inputs\_mu3tev\_v1"}.

\section{The Wilson coefficients, exactly --- and the load-bearing $1/6$}
\label{sec:codeWilsons}

In \code{compute\_delta\_f2\_wilsons} the coefficients are matched directly at
$\mu=\MKK$ with \code{prefactor = 1.0 / (couplings.M\_KK**2)} and four definite
numerical factors. The literal code lines (\code{deltaf2.py}:342--345) are
\begin{align}
\CVLL &= \texttt{left*left*prefactor/6.0}, &
\CVRR &= \texttt{right*right*prefactor/6.0},
\label{eq:codeC1}\\
\CLRfour &= \texttt{-(left*right)*prefactor}, &
\CLRfive &= \texttt{(left*right)*prefactor/3.0},
\label{eq:codeC45}
\end{align}
i.e.
\begin{align}
\CVLL &= \frac{(g_L)_{ij}^2}{6\,\MKK^2}, &
\CVRR &= \frac{(g_R)_{ij}^2}{6\,\MKK^2}, &
\CLRfour &= -\frac{(g_L)_{ij}(g_R)_{ij}}{\MKK^2}, &
\CLRfive &= \frac{(g_L)_{ij}(g_R)_{ij}}{3\,\MKK^2},
\label{eq:codeWilson}
\end{align}
where \code{left} $=(g_L)_{ij}$ and \code{right} $=(g_R)_{ij}$ are the complex
mass-basis KK-gluon couplings picked off by \code{\_pair\_couplings} for the
relevant pair (down-sector indices $(0,1)$ for the $s$--$d$ kaon entry, $(0,2)$
for $b$--$d$, $(1,2)$ for $b$--$s$; up-sector $(0,1)$ for the $c$--$u$ charm
entry). The crucial $1/6$ on the vector coefficient is \emph{not} a typo: it is
$\half$ (from the $t$-channel propagator structure) times $\tfrac13$ (from the
$SU(3)$ colour Fierz rearrangement of \eqref{eq:fierz}). The whole of
\S\ref{sec:r5} is about why this $1/6$ must never be ``corrected.''

\section{The hadronic matrix elements, exactly}
\label{sec:codeME}

The matrix elements live in two helpers. For the kaon,
\code{\_kaon\_matrix\_elements} returns the real $\langle O_i\rangle$ in the
code's $O_1/O_4/O_5$ basis; the generic
\code{\_meson\_matrix\_elements(f\_P, m\_P, m\_q1, m\_q2, B\_1, B\_4, B\_5)} does
the same for any pseudoscalar. Both build the same structure, with
$f_P^2 m_P$ the dimensionful prefactor and $r_\chi$ of \eqref{eq:rchi}:
\begin{align}
\langle O_1^{\mathrm{VLL}}\rangle &= \tfrac{1}{3}\,f_P^2 m_P\, B_1,
\label{eq:O1ME}\\
\langle O_4^{\mathrm{LR}}\rangle &= \big(\tfrac14\,r_\chi + \tfrac{1}{24}\big)\,f_P^2 m_P\, B_4,
\label{eq:O4ME}\\
\langle O_5^{\mathrm{LR}}\rangle &= \big(\tfrac{1}{12}\,r_\chi + \tfrac18\big)\,f_P^2 m_P\, B_5.
\label{eq:O5ME}
\end{align}
The code lines are \code{me\_vll = (1.0/3.0)*fp2\_mp*B\_1},
\code{me\_lr4 = (r\_chi/4.0 + 1.0/24.0)*fp2\_mp*B\_4}, and
\code{me\_lr5 = (r\_chi/12.0 + 1.0/8.0)*fp2\_mp*B\_5}
(\code{deltaf2.py}:1451--1453; the kaon copy at lines~1230--1232). The chiral factor
$r_\chi$ lives \emph{inside} the LR matrix elements \eqref{eq:O4ME}--\eqref{eq:O5ME}
--- the enhancement of \S\ref{sec:chiral}. The hadronic inputs are FLAG-2024 / PDG
constants pinned in the module; the per-system bag parameters are listed below.
\begin{center}
\begin{tabular}{lcccccc}
\toprule
System & $f_P$ (GeV) & $m_P$ (GeV) & $B_1$ & $B_4$ & $B_5$ & code constants \\
\midrule
$K$   & $0.1557$ & $0.49761$ & $0.5503$ & $0.903$ & $0.691$ & \code{F\_K},\,\code{M\_K},\,\code{B\_1\_K},\dots \\
$B_d$ & $0.1900$ & $5.27972$ & $0.87$ & $1.02$ & $0.96$ & \code{F\_BD},\,\code{M\_BD},\,\code{B\_1\_BD},\dots \\
$B_s$ & $0.2303$ & $5.36692$ & $0.87$ & $1.02$ & $0.96$ & \code{F\_BS},\,\code{M\_BS},\,\code{B\_1\_BS},\dots \\
$D$   & $0.2120$ & $1.86484$ & $0.75$ & $1.0$ & $1.0$ & \code{F\_D},\,\code{M\_D0},\,\code{B\_1\_D},\dots \\
\bottomrule
\end{tabular}
\end{center}
The $\tfrac13$ prefactor in \eqref{eq:O1ME} is the post-B3 M12-ready convention:
the usual $1/(2m_P)$ has already been absorbed. The $1/6$ in the Wilson
\eqref{eq:codeC1} is unchanged. See \S\ref{sec:r5}.

\section{Assembling $M_{12}^{\rm NP}$}
\label{sec:codeM12}

The generic assembler \code{compute\_m12\_np} forms \eqref{eq:m12sum} directly,
\begin{equation}
\Mtwelve^{\rm NP} = \CVLL\langle O_1\rangle + \CVRR\langle O_1\rangle
+ \CLRfour\langle O_4\rangle + \CLRfive\langle O_5\rangle,
\label{eq:codeM12}
\end{equation}
using the \emph{same} $\langle O_1\rangle$ for VRR as for VLL (parity), as the
code comment notes. The kaon has its own copy \code{\_compute\_m12\_np} (same
structure, kaon-specific matrix elements). From here the four observables
specialise:
\begin{itemize}[leftmargin=1.4em]
\item \textbf{$\eps_K$} (\code{evaluate\_epsilon\_k}): takes \emph{only}
\code{m12\_np.imag} and forms
$|\eps_K^{\rm NP}| = |\kappa_\eps/(\sqrt2\,\Delta m_K)\cdot\ImMtwelve|$ with
\code{KAPPA\_EPSILON = 0.94} and \code{DELTA\_M\_K = 3.484e-15}~GeV; the budget is
the SM-vs-experiment residual
$|\eps_K^{\rm exp}-\eps_K^{\rm SM}| = |2.228-2.161|\times10^{-3}=6.7\times10^{-5}$
(from \code{EPSILON\_K\_EXP} and \code{EPSILON\_K\_SM}, overridable). It is the
only observable that uses the imaginary part.
\item \textbf{$\Delta m_K$} (\code{evaluate\_delta\_mk}): the weaker companion,
$|\Mtwelve^{\rm NP}|<\Delta m_K/2$ (the magnitude, not the phase).
\item \textbf{$\Delta m_d$, $\Delta m_s$, $D^0$}
(\code{evaluate\_bd\_mixing}, \code{evaluate\_bs\_mixing},
\code{evaluate\_d0\_mixing}): each forms $|\Mtwelve^{\rm NP}|$ via
\code{compute\_m12\_np} with the system's hadronic inputs and compares to a budget.
The $B$ budgets are
\code{max($\Delta m^{\rm exp}/2$, $|\Delta m^{\rm exp}-\Delta m^{\rm SM}|/2$)}
(\code{\_bd\_budget}, \code{\_bs\_budget}); the $D$ budget is
$\Delta m_D^{\rm exp}/2$ alone (\code{\_d0\_budget}), because the SM long-distance
prediction is too dirty to subtract.
\end{itemize}
Each returns a result dataclass with \code{abs\_m12\_np}/\code{epsilon\_k\_np},
\code{budget}, \code{ratio\_to\_budget}, and \code{passes = ratio <= 1.0}.

\section{QCD running, exactly}
\label{sec:codeRunning}

By default \code{evaluate\_delta\_f2\_constraints} runs the Wilson coefficients
from the matching scale ($\MKK$) down to \code{mu\_had = 2.0}~GeV
(\code{apply\_qcd\_running=True}) via \code{\_evolve\_wilsons}, which calls
\code{qcd\_running.evolve\_deltaf2\_wilsons}. The running is leading-log, in the
Buras--Misiak--Urban basis $\{O_1^{\rm VLL}, O_1^{\rm VRR}, O_4^{\rm LR},
O_5^{\rm LR}\}$:
\begin{itemize}[leftmargin=1.4em]
\item the VLL/VRR operators run multiplicatively with a single anomalous dimension
$\gamma_{\rm VLL}=6(\Nc-1)/\Nc=4$ for $\Nc=3$ (\code{\_GAMMA\_VLL = 4.0}), so they
change only modestly;
\item the LR pair mixes through the $2\times2$ matrix \code{\_GAMMA\_LR}
$=\big[\!\begin{smallmatrix}-16 & -6\\ 0 & 2\end{smallmatrix}\!\big]$ (the
conjugated scalar-basis map of the BMU LR anomalous dimensions), which
\emph{enhances} the LR contribution by the $\sim8$ of \S\ref{sec:running}.
\end{itemize}
The strong coupling itself is evolved by \code{run\_alpha\_s} with
$\alpha_s(M_Z)=0.1179$ (\code{\_ALPHA\_S\_MZ\_DEFAULT}) and proper flavour
thresholds. There is a known sub-leading caveat (\S\ref{sec:bagscale}) that the LR
bag parameters are FLAG $3$~GeV values while the coefficients are run to $2$~GeV.

\section{Severity, the veto, and the per-system bounds}
\label{sec:codeVeto}

Each meson constraint is \code{Severity.HARD}: a point whose
\code{ratio\_to\_budget} exceeds $1$ is \emph{vetoed} (removed from the allowed
region). In the live hadronic path, the ratio is formed from the computed
observable itself: $|\varepsilon_K^{\rm NP}|$ for \code{K001} and
$|\Mtwelve^{\rm NP}|$ for the $B_d$, $B_s$, and $D$ magnitude constraints (with
the catalog wrappers overriding the core budget where needed). The
\code{dominant\_operator} and per-operator sizes remain diagnostics, useful for
seeing that LR terms usually dominate; they do \emph{not} replace the live
hadronic observable. Only the legacy \code{use\_hadronic=False} fallback uses the
dominant weighted operator as its conservative surrogate. The default experimental
inputs are pinned in \code{DEFAULT\_DELTA\_F2\_INPUTS\_V1}, with per-system
\code{bound} values (e.g.\ kaon \code{2.0e-8}, $B_d$ \code{1.667e-13}, $B_s$
\code{5.844e-12}, $D$ \code{3.125e-15}) and the LR-weighting flags
\code{lr1\_weight=7.0}, \code{lr2\_weight=2.0} that mark the LR dominance in the
legacy operator-weight path; the live HARD veto uses the hadronic path with the
budgets of \S\ref{sec:codeM12}.

% ====================================================================
\part{Subtleties and the gap to ``rigorous''}
\label{part:subtle}
% ====================================================================

\section{The normalization convention: the factor-of-two that nearly broke it}
\label{sec:r5}

This is the single most important subtlety in the whole pipeline. The R5 ledger
fixed the $1/6$ Wilson colour split. The later B3 audit fixed a separate
M12-normalisation issue in the matrix elements. The current convention is:
\begin{itemize}[leftmargin=1.4em]
\item \textbf{The repo convention} carries the colour factor in the
\emph{Wilson}: the vector coefficient gets a $1/6$ \eqref{eq:codeC1}, namely
$\half$ (propagator) $\times\,\tfrac13$ (colour Fierz, from \eqref{eq:fierz}).
\item \textbf{The repo matrix elements} are M12-ready GeV$^3$ objects:
$\langle O_1\rangle=\tfrac13 f_P^2 m_P B_1$,
$\langle O_4\rangle=(\tfrac14 r_\chi+\tfrac{1}{24})f_P^2 m_P B_4$, and
$\langle O_5\rangle=(\tfrac{1}{12} r_\chi+\tfrac18)f_P^2 m_P B_5$.
\item \textbf{The Lane-C \code{plpr\_projectors.v2} exports} store non-M12-ready
GeV$^4$ projector forms and divide by $2m_P$ when constructing $M_{12}$:
$Q_1=(2/3)f_P^2m_P^2B_1$, $Q_4=(1/2)r_\chi f_P^2m_P^2B_4$, and
$Q_5=(1/6)r_\chi f_P^2m_P^2B_5$ in the Lane-C file.
\end{itemize}
The pre-B3 review text collapsed these layers and certified the stale $\tfrac23$
M12-ready value. That statement is superseded.

\begin{callout}{Post-B3 correction: do not restore the pre-B3 matrix elements.}
The June-2026 B3 audit and the 2026-07 full-repo audit corrected this section.
The $1/6$ Wilson factor stays. The M12-ready matrix elements are the GGMS forms
in \eqref{eq:O1ME}--\eqref{eq:O5ME}. Do not restore the old $\tfrac23$ O1
coefficient or the swapped and doubled LR pair. The constraints
\code{K001}/\code{K002}/\code{B001}/\code{B003}/\code{C001} and the $\Delta F=2$
core are all in the post-B3 repo convention.
\end{callout}

\section{The $M_{\rm KK}$ convention: the factor of $2.45$}
\label{sec:mkkconv}

``$\MKK$'' is overloaded, and the overload propagates through the $1/\MKK^2$
prefactor in \eqref{eq:codeWilson}. The repo bookkeeping default is $\MKK=\LIR$
(\code{DEFAULT\_QUARK\_XI\_KK = 1.0}), but the \emph{physical} first KK gauge mass
is $x_1\LIR$ with the NN-boundary-condition root
\code{GAUGE\_KK\_ROOT\_NN = 2.448687135269161}. A bound stated as ``$20$~TeV'' in
$\LIR$ units is ``$\sim49$~TeV'' in physical first-KK units; conversely the
literature's ``$\sim21$~TeV'' is a \emph{physical} KK-gluon mass. Whichever
convention a quoted bound uses must match the $\MKK$ fed into
\code{compute\_delta\_f2\_wilsons}, where the $1/\MKK^2$ lives. Mixing conventions
here mis-scales the entire amplitude by $2.45^2\approx6$ --- a second, independent
factor hazard alongside the colour one of \S\ref{sec:r5}.

\section{The coupling model is an MFV fit, not a full 5D spectral computation}
\label{sec:couplingmodel}

The complex couplings $(g_L)_{ij},(g_R)_{ij}$ fed into \eqref{eq:codeWilson} come
from the repo's quark MFV-fit machinery
(\code{quarkConstraints.couplings.compute\_quark\_kk\_gluon\_couplings}), which
rotates the zero-mode overlap profiles into the exact quark mass basis returned by
the fit --- \emph{not} from a tower-summed five-dimensional mode computation. By
default it uses the perturbative $g_s(\MKK)$ unless a volume-enhanced $\gss$ is
supplied (\code{g\_s\_star}). This is the honest boundary of ``rigorous'': the
\emph{matching colour structure, hadronic matrix elements, RG running, and
CP-extraction} chain is rigorous and literature-cross-checked (\S\ref{sec:r5}),
but the \emph{coupling inputs} are a model layer. A fully first-principles
$\Delta F=2$ would replace them with explicit overlap integrals over the warped
profiles.

\section{Bag-parameter scale matching ($B_4,B_5$)}
\label{sec:bagscale}

A documented caveat lives in the core itself: the LR kaon bag inputs
\code{B\_4\_K\_3GEV}, \code{B\_5\_K\_3GEV} are FLAG-2024 $\overline{\rm MS}(3$~GeV)
values, while \code{mu\_had} defaults to $2$~GeV and the Wilson coefficients are
run to $2$~GeV. The matrix-element path keeps the $3$~GeV bag numbers for now so
evaluations stay numerically stable, with a \code{TODO} to RG-convert
$B_4,B_5$ to $2$~GeV in the same scheme. This is a small, known, sub-leading
mismatch --- flagged, not hidden.

\section{What is leading-log only / deferred / out of scope}
\label{sec:deferred}

To be complete about the gaps:
\begin{itemize}[leftmargin=1.4em]
\item \textbf{Leading-log running only.} \code{qcd\_running.py} implements
leading-log evolution; NLO matching at $\MKK$ and NLO running are not included.
\item \textbf{Dominant operators only.} The basis is
$\{\QVLL,\QVRR,\QLRfour,\QLRfive\}$; the additional SLL/SRR scalar operators are
not separately carried (sub-leading for these observables in this setup).
\item \textbf{Magnitude for the heavy mesons, phase only for the kaon.} $\eps_K$
uses $\ImMtwelve$; $\Delta m_{d,s}$ and $D$ use $|\Mtwelve|$. The $B_s$ phase
$\phi_s$ is handled by the companion constraint \code{B004} (argument of
$\Mtwelve$), not in this magnitude path.
\item \textbf{No flavour-symmetry model.} The pipeline takes whatever couplings the
fit produces; it does not itself impose alignment or a family symmetry (those would
be a different model input, \S\ref{sec:custodial}).
\end{itemize}

\section{Checklist: what a maximally rigorous $\Delta F=2$ would add}
\begin{enumerate}[leftmargin=1.8em]
\item Replace the MFV-fit couplings with explicit warped overlap integrals for
$(g_L)_{ij},(g_R)_{ij}$ in a fixed, documented $\MKK$ convention
(\S\ref{sec:couplingmodel}, \S\ref{sec:mkkconv}).
\item RG-convert $B_4,B_5$ from $3$~GeV to the $2$~GeV matching scheme
(\S\ref{sec:bagscale}).
\item Upgrade the leading-log running to NLO and add NLO matching at $\MKK$.
\item Carry the full SLL/SRR basis for completeness, even if sub-leading.
\item \textbf{Keep the R5 convention frozen} and guarded by the internal
normalisation tests and CFW comparison checks, so a future refactor cannot
silently reintroduce the factor-of-two
(\S\ref{sec:r5}). This is the highest-value guardrail in the list.
\end{enumerate}

% ====================================================================
\part*{Annotated reading guide}
\addcontentsline{toc}{section}{Annotated reading guide}
% ====================================================================

\begin{description}[leftmargin=0em,style=nextline]
\item[Csaki, Falkowski, Weiler, arXiv:0804.1954
(\emph{JHEP 0809:008}, 2008).] \emph{The} reference for the RS flavour problem and
the paper our $\Delta F=2$ core matches to machine precision (the R5 cross-check,
\S\ref{sec:r5}). The KK-gluon $\Delta F=2$ effective Hamiltonian (Eqs.~3.3--3.7),
the operator basis with the chirality-flipping $Q_4/Q_5$, the colour Fierz
decomposition $\sum_A\TA\otimes\TA=\half(\delta\delta-\tfrac1\Nc\delta\delta)$, the
chiral plus running enhancement of the LR operator, the much tighter constraint on
the \emph{imaginary} part of $C_4^{\rm LR}$ (their $\sim1.6\times10^5$~TeV vs.\
$\sim1.2\times10^4$~TeV on the coefficient), and the headline $M_G\gtrsim22\pm6$~TeV
(generic anarchic) / $\sim33$~TeV (pseudo-Goldstone) KK-gluon floor. Read
\S\ref{sec:treelevel}--\S\ref{sec:unified} alongside it.

\item[Blanke, Buras, Duling, Gori, Weiler, arXiv:0809.1073.] $\Delta F=2$ in the
\emph{custodial} RS model. Their Eq.~4.33 is the meson $\Mtwelve$ from KK-gluon
exchange in the $Q_1^{\rm VLL},Q_4^{\rm LR},Q_5^{\rm LR}$ basis; they state
explicitly that custodial $P_{LR}$ protects the $Z\bar d_L d_L$ vertex but does
\emph{not} protect $\Delta F=2$, and that $\eps_K$ still forces
$\MKK\gtrsim20$~TeV. The reference for the ``custodial does not fix this'' point of
\S\ref{sec:custodial}, and the second leg of the R5 cross-check (ratio $\sim1.69$,
a basis/running difference, not the failure mode).

\item[Agashe, Perez, Soni, arXiv:hep-ph/0408134.] The flavour structure of warped
models and the RS-GIM mechanism (\S\ref{sec:rsgim}). Establishes that anarchic 5D
Yukawas fix the flavour parameters through the zero-mode profiles, and that
$\eps_K$ from the misaligned doublet bulk-mass dispersion is the sharpest bound,
pushing the cutoff well above the few-TeV scale wanted for naturalness.

\item[Buras, Misiak, Urban, \emph{Nucl.\ Phys.\ B} (2000).] The QCD running of the
$\Delta F=2$ operators: the BMU operator basis and the anomalous dimensions our
\code{qcd\_running.py} encodes (\code{\_GAMMA\_VLL = 4.0} and the $2\times2$
\code{\_GAMMA\_LR}; \S\ref{sec:codeRunning}). The reference for the $\eta_1^{-5}\sim8$
LR running enhancement of \S\ref{sec:running}.

\item[FLAG (Flavour Lattice Averaging Group), 2024 review.] The source of the
hadronic matrix-element inputs --- decay constants $f_P$ and bag parameters
$B_1,B_4,B_5$ --- pinned in \code{deltaf2.py} and tabulated in \S\ref{sec:codeME}.
The chiral-factor light-quark masses ($m_s,m_d$ at $2$~GeV) are also FLAG.

\item[The four sibling notes
(\code{epsilon\_k\_review}, \code{d0\_mixing\_review},
\code{delta\_m\_d\_review}, \code{delta\_m\_s\_review}).] Each specialises this
framework to one observable: the kaon note develops $\eps_K$ and the imaginary
part in detail; the charm note the long-distance budget; the two beauty notes the
magnitude (and, for $B_s$, the phase $\phi_s$ via \code{B004}). All four inherit
the operator basis, the matrix elements, and --- crucially --- the R5
normalization convention of this note. \textbf{When a meson note and this note seem
to disagree on a factor, this note is canonical (\S\ref{sec:r5}).}
\end{description}

\bigskip
\noindent\rule{\linewidth}{0.4pt}\\
\footnotesize
\emph{Internal note. The operator basis \eqref{eq:QVLL}--\eqref{eq:QLR5}, the
matching \eqref{eq:codeWilson}, and the matrix elements
\eqref{eq:O1ME}--\eqref{eq:O5ME} are the exact forms our code uses; the numerical
constants ($1/6$ Wilson factor, $\tfrac13$ M12-ready $O_1$, $r_\chi$, the bag parameters,
$\kappa_\eps=0.94$, $\Delta m_K=3.484\times10^{-15}$~GeV) are as stored in
\code{quarkConstraints/deltaf2.py}. The ``rigorous vs.\ proxy'' tags and the R5
resolution follow the orchestration ledger
\code{.orchestration/PHASE2\_PROGRAM\_LEDGER.md}. Cross-check before citing
externally; this is the authoritative normalization reference for the four meson
notes.}

\end{document}
